\documentclass[11pt]{article}
\usepackage[margin=1in]{geometry}
\usepackage{amsmath,amssymb,amsthm,bm,mathtools}
\usepackage[numbers]{natbib}
\usepackage[colorlinks=true,linkcolor=blue,citecolor=blue]{hyperref}

\newcommand{\N}{\mathcal{N}}
\newcommand{\E}{\mathbb{E}}
\newcommand{\R}{\mathbb{R}}
\newcommand{\KL}{\mathrm{KL}}
\newcommand{\BF}{\mathrm{BF}}
\newcommand{\simplex}{\Delta}
\DeclareMathOperator*{\argmax}{arg\,max}
\DeclareMathOperator*{\argmin}{arg\,min}
\newtheorem{proposition}{Proposition}
\newtheorem{lemma}{Lemma}

\title{\bf An implicit-weight mixture-effect SER:\\ clean derivation}
\author{MixSER notes}
\date{\today}

\begin{document}
\maketitle

\noindent
This note derives, from scratch, the mixture-effect single-effect regression
(SER) used in \texttt{MixSER}, with the mixing weights learned \emph{implicitly}
by empirical Bayes -- the device behind \texttt{ashr} \citep{stephens2017},
solved by the convex optimiser \texttt{mixsqp} \citep{kim2020} -- rather than
through a Dirichlet prior. There are two layers: a \emph{coding-model} mixture at
the SNP level (\S\ref{sec:level1}, identifiable, the recommended core) and a
\emph{per-individual} mixture within a SNP (\S\ref{sec:level2}). Both reduce to
the same convex program (\S\ref{sec:eb}). Throughout, $\sigma^2$ is the residual
variance and effects are in residual-SD units.

\section{The single-effect building block}
\label{sec:ser}
Let $x\in\R^n$ be one centred (and scaled) predictor and $y\in\R^n$ a centred
response. The Bayesian simple regression
\begin{equation}
y \mid \beta \sim \N(x\beta,\sigma^2 I_n),\qquad \beta \sim \N(0,\sigma_0^2)
\label{eq:simple}
\end{equation}
has a closed-form marginal. With $\hat\beta = x^\top y/x^\top x$ and
$s^2 = \sigma^2/x^\top x$, completing the square in $\beta$ gives
$p(y\mid\sigma_0^2)=\N(y;0,\sigma^2 I + \sigma_0^2 x x^\top)$, and the log Bayes
factor against $\sigma_0^2=0$ is the Wang--Sarkar--Carbonetto--Stephens form
\citep{wang2020},
\begin{equation}
\log\BF(\sigma_0^2)
= \tfrac12\log\frac{s^2}{\sigma_0^2+s^2}
   + \tfrac12\,\hat\beta^2\Big(\frac{1}{s^2}-\frac{1}{\sigma_0^2+s^2}\Big),
\label{eq:lbf}
\end{equation}
with Gaussian posterior
\begin{equation}
\beta\mid y \sim \N(\mu,\tau^2),\qquad
\tau^2=\big(\sigma_0^{-2}+s^{-2}\big)^{-1},\qquad
\mu=\tfrac{\tau^2}{s^2}\hat\beta .
\label{eq:postbeta}
\end{equation}
The single-effect regression \citep{wang2020} places one effect among $p$
candidates with prior weights $\gamma_j$ (uniform $1/p$): posterior inclusion
$\alpha_j \propto \gamma_j\BF_j$, SER log Bayes factor $\log\sum_j\gamma_j\BF_j$,
and prior variance by empirical Bayes
$\hat\sigma_0^2=\argmax_{\sigma_0^2}\log\sum_j\gamma_j\BF_j(\sigma_0^2)$.

\section{Empirical-Bayes mixture weights (the \texttt{ash}/\texttt{mixsqp} core)}
\label{sec:eb}
Suppose we have $T$ ``observations'' $t$ and $K$ components, with likelihood
$L_{tk}\ge0$ of observation $t$ under component $k$, and weights
$w\in\simplex^{K-1}=\{w\ge0,\mathbf 1^\top w=1\}$. The marginal log-likelihood
\begin{equation}
\ell(w)=\sum_{t=1}^T \log\Big(\sum_{k=1}^K w_k L_{tk}\Big)
\label{eq:ashobj}
\end{equation}
is maximised over the simplex.

\begin{proposition}[Convexity]
\label{prop:convex}
$\ell(w)$ is concave on $\simplex$, so $\hat w=\argmax_{w\in\simplex}\ell(w)$ is a
convex program.
\end{proposition}
\begin{proof}
Each $w\mapsto L_{t\cdot}w$ is linear; $\log$ is concave and nondecreasing; the
composition is concave, a sum of concave functions is concave, and $\simplex$ is
convex.
\end{proof}

\noindent This is the maximum-likelihood mixture-proportions problem of
\texttt{ash} \citep{stephens2017}, solved by SQP in \texttt{mixsqp}
\citep{kim2020}. Two facts we reuse:

\begin{lemma}[Row scaling is irrelevant]
\label{lem:scale}
Replacing $L_{tk}\leftarrow c_t L_{tk}$ ($c_t>0$) leaves $\argmax\ell$ unchanged
(it shifts $\ell$ by $\sum_t\log c_t$). Hence we may feed \emph{Bayes factors}.
\end{lemma}

\begin{lemma}[EM is an MM algorithm for \eqref{eq:ashobj}]
\label{lem:em}
With $r_{tk}= w_k L_{tk}/\sum_{k'}w_{k'}L_{tk'}$, the update
$w_k \leftarrow \tfrac1T\sum_t r_{tk}$ increases $\ell$ monotonically to $\hat w$.
\end{lemma}
\begin{proof}
By Jensen $\ell(w')-\ell(w)\ge\sum_{t,k}r_{tk}\log(w'_k/w_k)$, maximised over
$\simplex$ at $w'_k=\tfrac1T\sum_t r_{tk}$; the MM principle gives monotonicity
and convergence to the convex optimum.
\end{proof}

\noindent The implementation uses \texttt{mixsqp} when available and
Lemma~\ref{lem:em} as a fallback. The responsibility $r_{tk}$ is the posterior
weight that observation $t$ belongs to component $k$.

\section{Level 1: the coding-model mixture SER}
\label{sec:level1}
A genotype $g\in\{0,1,2\}$ gives codings $x^A(g)=g$, $x^D(g)=\mathbf1\{g\ge1\}$,
$x^R(g)=\mathbf1\{g=2\}$. Inside one IBSS single effect we hold a centred partial
residual $y$ and, for variant $j$ and coding $m$, the centred-scaled $x_j^m$.

\paragraph{Per-coding hyperparameters (as in SuSiE).}
\begin{equation}
\hat\sigma_{0m}^2=\argmax_{V}\ \log\sum_{j}\gamma_j\,\BF_{jm}(V),
\qquad \BF_{jm}(V)=\exp(\text{\eqref{eq:lbf} for }x_j^m),
\label{eq:level1-V}
\end{equation}
allowing $\hat\sigma_{0m}^2=0$ (coding off) when that does not decrease the
objective -- this is what makes the procedure additive-leaning without tuning.

\paragraph{Implicit model weights (this replaces the Dirichlet).}
Augment with an explicit null ($\BF_{j0}\equiv1$) of fixed prior mass
$w_\varnothing$ (default $0$, as \texttt{susieR}'s \texttt{null\_weight}). The
coding weights $w=(w_A,w_D,w_R)\in\simplex^2$ are learned by \S\ref{sec:eb} with
the $p$ variants as observations and codings as components,
\begin{equation}
\hat w=\argmax_{w\in\simplex^2}\ \sum_{j=1}^{p}\log\Big(\sum_{m\in\{A,D,R\}}
w_m\,\BF_{jm}(\hat\sigma_{0m}^2)\Big),
\label{eq:level1-w}
\end{equation}
convex by Prop.~\ref{prop:convex}, valid with Bayes factors by
Lemma~\ref{lem:scale}. It is self-regularising: a null variant has
$\BF_{jm}\approx1$, contributing $\log(\sum_m w_m)=0$, so only signal variants
move $\hat w$. (Genome-wide: solve \eqref{eq:level1-w} once over many SNPs and
freeze $\hat w$ -- the pooled mode.)

\paragraph{Posteriors.} By Bayes' rule, for $m\in\{A,D,R\}$,
\begin{equation}
\alpha_{jm}
=\frac{\gamma_j(1-w_\varnothing)\hat w_m\BF_{jm}}
{\,w_\varnothing+\sum_{j'm'}\gamma_{j'}(1-w_\varnothing)\hat w_{m'}\BF_{j'm'}\,},
\label{eq:level1-post}
\end{equation}
giving PIP $\alpha_j=\sum_m\alpha_{jm}$, coding posterior
$\Pr(m\mid j)=\alpha_{jm}/\alpha_j$, SER log-BF
$\log\sum_j\gamma_j\sum_m\hat w_m\BF_{jm}$, and fitted effect
$\hat y=\sum_{j,m}\alpha_{jm}x_j^m\mu_{jm}$ -- exactly the inputs IBSS consumes.

\begin{proposition}[Exact reduction to SuSiE]
With one coding and $w_\varnothing=0$, \eqref{eq:level1-post} is
$\alpha_j=\gamma_j\BF_j/\sum_{j'}\gamma_{j'}\BF_{j'}$ with $\hat\sigma_0^2$ from
\eqref{eq:level1-V}: the standard additive SuSiE SER, identically (verified:
$\max_j|\alpha_j^{\mathrm{mix}}-\alpha_j^{\mathrm{SuSiE}}|=0$).
\end{proposition}

\noindent Level 1 is a statement about the conditional means $\E[y\mid g]$,
identifiable at $O(\theta^2)$; it is where every coding decision should be made.

\section{Level 2: per-individual mixture, without a Dirichlet}
\label{sec:level2}
At a single variant, each individual $i$ carries a latent coding
$z_i\in\{1,\dots,M\}$ with fractions $\pi\in\simplex^{M-1}$:
\begin{equation}
z_i\sim\mathrm{Cat}(\pi),\quad
y_i\mid z_i{=}m \sim \N(x_i^m\beta_m,\sigma^2),\quad
\beta_m\sim\N(0,\sigma_{0m}^2).
\label{eq:pim}
\end{equation}
Treat $\pi,\sigma_{0m}^2,\sigma^2$ as hyperparameters estimated by empirical
Bayes -- so there is \textbf{no} $q(\pi)$ and \textbf{no} Dirichlet KL. With the
mean-field posterior $\prod_i q(z_i)\prod_m q(\beta_m)$, $q(z_i{=}m)=r_{im}$,
$q(\beta_m)=\N(\mu_m,\tau_m^2)$, the evidence lower bound is
\begin{equation}
\begin{aligned}
\mathcal F=\;&
-\tfrac{n}{2}\log(2\pi\sigma^2)-\tfrac{1}{2\sigma^2}\sum_i y_i^2
+\sum_i\sum_m r_{im}\Big[\log\pi_m
-\tfrac{1}{2\sigma^2}\big({-}2x_i^m y_i\mu_m+(x_i^m)^2(\mu_m^2+\tau_m^2)\big)\Big]\\
&-\sum_i\sum_m r_{im}\log r_{im}
-\sum_m \KL\!\big(\N(\mu_m,\tau_m^2)\,\|\,\N(0,\sigma_{0m}^2)\big),
\end{aligned}
\label{eq:elbo}
\end{equation}
with $\KL=\tfrac12(\log\tfrac{\sigma_{0m}^2}{\tau_m^2}
+\tfrac{\tau_m^2+\mu_m^2}{\sigma_{0m}^2}-1)$. The two data constants are kept so
$\mathcal F$ is a genuine lower bound on $\log p(y)$.

\paragraph{Coordinate ascent.}
\begin{align}
q(\beta_m):\ &
\tau_m^2=\Big(\sigma_{0m}^{-2}+\sigma^{-2}\textstyle\sum_i r_{im}(x_i^m)^2\Big)^{-1},
\ \ \mu_m=\tfrac{\tau_m^2}{\sigma^2}\textstyle\sum_i r_{im}x_i^m y_i,
\label{eq:u-beta}\\
q(z_i):\ &
\log r_{im}\propto \log\pi_m
-\tfrac{1}{2\sigma^2}\big({-}2x_i^m y_i\mu_m+(x_i^m)^2(\mu_m^2+\tau_m^2)\big),
\label{eq:u-z}\\
\sigma_{0m}^2:\ & \sigma_{0m}^2=\mu_m^2+\tau_m^2,
\label{eq:u-s0}\\
\pi:\ &
\pi=\argmax_{\pi\in\simplex}\sum_i\log\big(\textstyle\sum_m \pi_m L_{im}\big),
\ \ L_{im}=\exp(\E_q\log\N(y_i;x_i^m\beta_m,\sigma^2)),
\label{eq:u-pi}\\
\sigma^2:\ &
\sigma^2=\tfrac1n\sum_i\sum_m r_{im}\big(y_i^2-2x_i^m y_i\mu_m+(x_i^m)^2(\mu_m^2+\tau_m^2)\big).
\label{eq:u-s2}
\end{align}
Equation \eqref{eq:u-z} is the softmax of \eqref{eq:elbo}; \eqref{eq:u-s0} the EB
maximiser of the KL term; \eqref{eq:u-beta} the weighted \eqref{eq:postbeta}; and
\eqref{eq:u-s2} the stationary point $\partial\mathcal F/\partial\sigma^2=0$.
\textbf{Estimating $\sigma^2$ is essential for a standalone single-SNP fit}: with
$\sigma^2$ fixed too large, the responsibilities collapse to the prior $\pi$
(everyone looks like a blend); EB recovers the noise level and with it sharp
responsibilities. (Inside SuSiE, $\sigma^2$ is supplied and \eqref{eq:u-s2} is
skipped.)

\paragraph{Why the $\pi$-update \emph{is} the \texttt{mixsqp} program.}
The $\pi$-terms of \eqref{eq:elbo} are $\sum_{i,m}r_{im}\log\pi_m$; at optimum
$r_{im}\propto\pi_m L_{im}$ \eqref{eq:u-z}, so profiling $r$ out turns them into
$\sum_i\log\sum_m\pi_m L_{im}$ -- exactly \eqref{eq:ashobj} with individuals as
observations and codings as components. This is the Dirichlet's replacement:
where the old code formed $\E[\log\pi_m]$ from a Dirichlet posterior, we set
$\log\pi_m$ from the \texttt{mixsqp} solution; the EM form is
$\pi_m=\tfrac1n\sum_i r_{im}$.

\paragraph{Initialisation.} The uniform point $r_{im}\equiv1/M$ is a saddle of
\eqref{eq:elbo}, and random starts can fall into degenerate optima. Initialise
from the \emph{pure single-coding} least-squares effects
$\mu_m^{(0)}=\sum_i x_i^m y_i/\sum_i(x_i^m)^2$ and the responsibilities they
imply: this breaks symmetry in a data-driven direction and attains the highest
ELBO (a few random restarts are kept as insurance, selecting the best
$\mathcal F$).

\paragraph{Reading $r$.} $r_{im}$ is sharp (near $0/1$) only for individuals
whose codings genuinely differ -- at $g=2$, where additive and recessive
predictions diverge. At $g=0$ all codings predict the baseline, the data are
uninformative, and $r_{i\cdot}$ \emph{correctly} returns the prior $\pi$.
Summarise $r$ stratified by genotype, never pooled over $g$.

\section{Why decisions live at Level 1, and the pooling requirement}
\label{sec:decisions}
The per-individual fractions are weakly identified: with $g\in\{0,1,2\}$ the data
are summarised by $p(y\mid g)$, and the mixture's signal over a single-coding
mean is an excess within-genotype variance
$\mathrm{Var}(y\mid g{=}1)=\sigma^2+\pi(1-\pi)\theta^2$ -- $O(\theta^2)$ in a
second moment, $O(\theta^4)$ in information, confounded with heteroscedasticity
(full analysis in \texttt{identifiability\_per\_individual\_mixture.md}). Hence:

\begin{enumerate}
\item \textbf{Selection is a Level-1 task.} The
additive/dominant/recessive/mixture posterior is \eqref{eq:level1-post}, each
model scored by a proper marginal likelihood. One must not decide pure-vs-mixture
by a per-locus point estimate of $\pi$: \eqref{eq:pim} nests the pure models (put
$\pi$ on a vertex), so $\max_\pi p(y\mid\pi)\ge p_{\mathrm{pure}}(y)$ always --
the maximised-$\pi$ mixture wins by construction (the \texttt{VEB-pi} pathology;
integrating $\pi$ under a Dirichlet was the old fix).
\item \textbf{$\pi$ must be pooled across loci.} Within one locus even the convex
MLE \eqref{eq:u-pi} places interior mass to fit within-genotype noise. The honest
EB cure -- the one that makes \texttt{ash} valid -- estimates \emph{one} $\pi$
across many loci by stacking the $L_{im}$ of \eqref{eq:u-pi} into a single
\eqref{eq:ashobj}; a pure-additive locus then contributes under the shared $\pi$
and cannot manufacture a mixture. Per-individual $r_{im}$ are reported
conditionally on this pooled $\hat\pi$; per-locus $\pi$ is exploratory only.
\end{enumerate}

\noindent The Dirichlet's two roles -- an Occam factor and regularising $\pi$ --
are taken over by Level-1's pooled marginal likelihood and by cross-locus EB
pooling of \eqref{eq:u-pi}, with no prior to tune.

\section*{Defaults and what ships}
Codings $\{A,D,R\}$ ($+$ optional null, $w_\varnothing=0$); $\sigma_{0m}^2$ by
\eqref{eq:level1-V}; coding weights $\hat w$ by \eqref{eq:level1-w}
(\texttt{per\_ser} or pooled); $\pi$ by \eqref{eq:u-pi}; $\sigma^2$ by
\eqref{eq:u-s2}; solver \texttt{mixsqp} with EM fallback. Implemented in
\texttt{ser\_mix\_eb()} (Level 1), \texttt{ser\_pim\_vb\_eb()} (Level 2,
pure-coding init, EB $\sigma^2$), and \texttt{estimate\_mixture\_weights()}.

\begin{thebibliography}{9}
\bibitem{stephens2017} M.~Stephens (2017). False discovery rates: a new deal.
\emph{Biostatistics} 18(2):275--294.
\bibitem{kim2020} Y.~Kim, P.~Carbonetto, M.~Stephens, M.~Anitescu (2020). A fast
algorithm for maximum likelihood estimation of mixture proportions using
sequential quadratic programming. \emph{JCGS} 29(2):261--273.
\bibitem{wang2020} G.~Wang, A.~Sarkar, P.~Carbonetto, M.~Stephens (2020). A simple
new approach to variable selection in regression, with application to genetic
fine mapping. \emph{JRSS-B} 82(5):1273--1300.
\end{thebibliography}

\end{document}
